\input{styles/article.tex}

% 公式排版
\usepackage{amsmath}
% 数学符号
\usepackage{amssymb}
% SI 制单位
\usepackage{siunitx}

\title{电的物理性质}
\date{2024-08-11}

\begin{document}
% 生成标题和目录
\maketitle\newpage
\tableofcontents\newpage

\section{电流}

\textbf{电流}（electric current）是电荷在电场或（半）导体内的平均定向移动。\textbf{电路的方向定义为正电荷移动的方向}，电路的大小则称为电流强度，是指单位时间内通过导线某一截面的电荷净转移量，每秒通过 1 库伦的电荷量称为 1 安培。“电流强度”也常直接简称为“电流”或称为“电流量”。
{\color{red}电流的方向被定义为与正电荷在电路中的移动方向相同，但实际上并不是正电荷移动，而是负电荷移动。}\textbf{电子流}是自由电子（不受原子束缚的负电子）在电路中的移动，其方向为电流的反向。

\section{电阻}

\subsection{欧姆定律}

电导体的两端电压 $U$ 与通过导体的电流 $I$ 成正比，表达式为：$U \propto I$，电阻 $R$ 与电压 $U$ 和电流 $I$ 的关系为：$R = \frac{U}{I}$，以此类推：$U = IR$。

\subsection{电阻率}

设导体的电阻为 $R$，电阻率为 $p$，横截面积为 $S$，长度为 $L$，它们的关系如下：

\begin{align*}
% &= 表示等号对齐
p &= R \times S \div L \\
R &= p\frac{L}{S}
\end{align*}

导体的电阻率一般会随着温度的变化而变化，可能会升高，也可能会减小。

以铜为例，在 \SI{20}{\degreeCelsius} 下，铜的电阻率为 \SI{1.7e-8}{\ohm/\square\metre}，那么长度为 \SI{100}{\square\metre}，横截面积为 \SI{0.8}{\square\milli\metre} 的铜线的电阻是多少？
\begin{align*}
R &= 1.7\times10^{-8}\times\frac{100}{0.8\times10^{-8}} \\
&= 0.000000017\times12500000000 \\
&= \SI{212.5}{\ohm}
\end{align*}

\section{电的其它物理量}
\subsection{电能量}
\SI{1}{\joule} 等于 \SI{1}{\watt/\second}，\SI{1}{\kWh} = \SI{3.6e6}{\joule}。

\subsection{电功率}
电功率 $P = U \times I$，其它相关公式如下：
\begin{align*}
&\because P = U \times I, U = I \times R \\
&\therefore P = (I \times R) \times I = I^2 \times R \\
&\because I = U \div R \\
&\therefore P = U \times (U \div R) = U^2 \div R
\end{align*}

\end{document}